\documentclass{situationist-preprint}

\title{Assessing Actual Markets}
\msid{POL-12-b}
\msfamily{POL-12}
\msdiscipline{Political Philosophy}
\mstype{Companion manuscript}
\mskeywords{markets; justice; difference principle; housing; credit; campaign finance}

\begin{document}
\maketitle

\begin{abstract}
POL-12-a argues that a market is just when its rules respect the basic liberties and no feasible change in those rules, the rest of the regime held fixed, yields an outcome better by the modified difference principle. It concludes that many, but not all, actual markets are just. This companion turns the criterion into a procedure for assessing particular markets and applies it to four: retail food markets, labor markets with collective bargaining, residential housing markets under exclusionary zoning, and markets in political advertising. The first two pass in the regimes we consider; the last two fail, for different reasons. The verdicts illustrate how the criterion separates the justice of markets from the voluntariness of exchange.
\end{abstract}

\section{A procedure}

By Definition~3.4 and Theorem~4.1 of POL-12-a, the justice of a market $m$ in a regime $R$ is settled by two questions, which we put as steps.

\begin{definition}\label{d:procedure}
To \emph{assess} a market $m$ in a regime $R$: (S1) determine whether the rules of $m$ respect equal basic liberty, including, as Proposition~5.5 of POL-12-a requires, whether they permit basic liberties to be exchanged or alienated; (S2) identify the feasible alternative rules for $m$, holding the other institutions of $R$ fixed, and ask whether any of them raises the expected level of the worst-off position or, at that level, the total.
\end{definition}

Step S2 is comparative and local. It does not ask whether $R$ is just; by Proposition~5.2 of POL-12-a a market may pass while its regime fails. Nor does it ask whether the exchanges are voluntary, which Theorem~4.1 of POL-12-a shows to be beside the point. Its difficulty is empirical: the effects of alternative rules on the worst-off are often uncertain, and an assessment is only as good as the evidence about them. Where the evidence does not discriminate, the verdict is that the market is not shown to be unjust.

\section{Four markets}

\emph{Retail food.} In a regime with an income floor delivered through transfers, retail food markets with ordinary rules of contract, safety and competition pass S1. For S2, the feasible alternatives are price controls and rationing. Both are known to produce shortages and queues, whose costs fall heavily on those with least time and least access to alternative supply \citep{Sen1981}; given the transfers, the worst-off do better under the market rules. The market is just.

\emph{Labor with collective bargaining.} Where workers may organize and bargain collectively, the rules of the labor market protect freedom of association and pass S1. The relevant alternatives under S2 are rules that weaken bargaining rights and rules that fix wages by administrative order. The first lower the share of the least advantaged workers without a compensating gain in their position through employment; the second sacrifice the information that wages carry about where labor is needed. Within the range of feasible rules, those that secure bargaining rights and a legal minimum do at least as well for the worst-off as any alternative we can identify. The market is just in regimes that have them, and in that respect unjust in regimes that do not.

\emph{Housing under exclusionary zoning.} Residential housing markets in which local rules prohibit all but detached houses on large lots pass S1 but fail S2. The rules restrict supply, raise rents and exclude lower-income households from places with better schools and employment, and the alternative rules that permit denser building are feasible with the rest of the regime fixed. The beneficiaries are existing owners, who are not the worst-off. That every sale and lease in such a market is voluntary does not affect the verdict.

\emph{Political advertising.} A market in which access to political speech is allocated by unlimited private expenditure fails S1. The equal political liberties require not only formal rights to vote and speak but their fair value, so that persons similarly motivated and endowed have roughly equal chances to influence political decisions whatever their economic position \citep[\S36]{Rawls1971}. A market that converts wealth into political influence without limit sells, in effect, a share in a basic liberty, and Proposition~5.5 of POL-12-a applies. Rules that cap expenditure or fund campaigns publicly remove the defect.

\begin{objection}
The verdicts depend on contestable empirical claims; a libertarian or a socialist would reach different ones.
\end{objection}

\begin{reply}
They would, and the criterion explains where the disagreement lies. Disputes about S2 are disputes about the effects of rules on the worst-off, which evidence can settle in principle; the libertarian's appeal to voluntariness and the socialist's appeal to the commodity form are, by Theorem~4.1 of POL-12-a, not the right kind of consideration. The criterion does not make assessment easy; it makes clear what an assessment must show.
\end{reply}

\section{Conclusion}

Applied to particular markets, the criterion of POL-12-a gives mixed verdicts: some markets, such as retail food and labor with collective bargaining in a regime with a floor, are just; others, such as exclusionary housing markets and unlimited markets in political speech, are not. This is what the conclusion of POL-12 asserts: many, but not all, actual markets are just. The case for or against a given market is decided by its rules and their effects on the worst-off, not by the voluntariness of its exchanges \citep{Satz2010}.

\begin{thebibliography}{9}\small
\bibitem[Rawls(1971)]{Rawls1971} Rawls, J. (1971). \emph{A Theory of Justice}. Harvard University Press.
\bibitem[Satz(2010)]{Satz2010} Satz, D. (2010). \emph{Why Some Things Should Not Be for Sale: The Moral Limits of Markets}. Oxford University Press.
\bibitem[Sen(1981)]{Sen1981} Sen, A. (1981). \emph{Poverty and Famines: An Essay on Entitlement and Deprivation}. Clarendon Press.
\end{thebibliography}

\end{document}
